\pdfoutput=1
\documentclass[12pt]{article}

\usepackage{amsmath, amssymb, amsthm}
\usepackage{graphicx}
\usepackage{booktabs}
\usepackage{array}
\usepackage{longtable}
\usepackage{calc}
\usepackage{etoolbox}
\makeatletter
\patchcmd\longtable{\par}{\if@noskipsec\mbox{}\fi\par}{}{}
\makeatother
\IfFileExists{footnotehyper.sty}{\usepackage{footnotehyper}}{\usepackage{footnote}}
\makesavenoteenv{longtable}
\usepackage[colorlinks=true, linkcolor=blue, citecolor=blue, urlcolor=blue]{hyperref}
\usepackage{geometry}
\geometry{margin=1in}

\makeatletter
\newsavebox\pandoc@box
\newcommand*\pandocbounded[1]{%
  \sbox\pandoc@box{#1}%
  \Gscale@div\@tempa{\textheight}{\dimexpr\ht\pandoc@box+\dp\pandoc@box\relax}%
  \Gscale@div\@tempb{\linewidth}{\wd\pandoc@box}%
  \ifdim\@tempb\p@<\@tempa\p@\let\@tempa\@tempb\fi
  \ifdim\@tempa\p@<\p@\scalebox{\@tempa}{\usebox\pandoc@box}%
  \else\usebox{\pandoc@box}%
  \fi
}
\makeatother
\providecommand{\tightlist}{%
  \setlength{\itemsep}{0pt}\setlength{\parskip}{0pt}}

\title{Power-of-Two Convergence Counts in the Tholonic Seed Space:\\
\large A Complete Hierarchy of Classical Constants}
\author{J. W. Milton\\[4pt]{\small Clarity Coalition}}
\date{\small 8 June 2026 \\ v1.1}

\begin{document}
\maketitle

\noindent\textbf{Keywords:} tholonic model, N-D-C triad, five constants, golden ratio, pi, Euler's number, natural logarithm, seed-space sweep, convergence hierarchy, generative entropy, power of two, Fibonacci seeds, archetype

\begin{abstract}
Paper 1 of this series establishes that five classical mathematical
constants (\(\pi/4\), \(\varphi\), \(e\), \(\sqrt{2}\), \(\ln 2\))
emerge as limits of a single family of three-variable tholonic
recurrences parameterized by an initial triple \((N_0, D_0, C_0)\) with
canonically distinct seeds. This paper extends that structural analysis
by an exhaustive computational sweep: all \(4^4 = 256\) combinations of
the four canonical seed values \(\{1, 2, 3, 5\}\) are applied to each of
the five algorithms, yielding 1,280 independent convergence runs.

The central finding is that the number of seed combinations converging
to each known constant is an exact power of two: \(\varphi\) and
\(\sqrt{2}\) each attract \(2^8 = 256\) combinations (every seed works),
\(\ln 2\) attracts \(2^6 = 64\), \(e\) attracts \(2^4 = 16\), and
\(\pi/4\) is reached by exactly \(2^0 = 1\) combination, the unique
Fibonacci triple \((N, I, D, C) = (1, 2, 3, 5)\). When the standard
\(\times 4\) scaling is applied to restore \(\pi\) from the raw Leibniz
output \(\pi/4\), the predicted slot of \(2^2 = 4\) is filled and the
sequence of even powers \(2^0, 2^2, 2^4, 2^6, 2^8\) is complete and
gapless.

We interpret this hierarchy through the lens of information-theoretic
entropy: \(\pi/4\) has zero generative entropy (one path, no choice),
while \(\varphi\) and \(\sqrt{2}\) have maximum generative entropy
(every path leads there). The cascade divides by exactly 4 at each step,
producing a four-stage emanation structure. We argue that this hierarchy
is not an artifact of the chosen seed set \(\{1, 2, 3, 5\}\) but
reflects the algebraic character of the constants themselves: their
positions in the hierarchy mirror their standing in the algebraic
hierarchy from quadratic irrationals (\(\varphi\), \(\sqrt{2}\)) through
transcendentals (\(e\), \(\ln 2\), \(\pi\)).

The paper does not introduce new proofs of convergence for the
individual series; those are classical and established in Paper 1. The
contribution is the combinatorial structure of the seed space, the
entropy interpretation, and the identification of \(\pi\) as an
archetype rather than an instance within this framework.
\end{abstract}

\section{Introduction}\label{introduction}

\subsection{Background}\label{background}

Paper 1 of this series {[}Mil26a{]} establishes that five classical
constants, \(\pi/4\) (via the tholonic Class A advancing branch, which
reduces mathematically to a grouping of the Leibniz series), \(\varphi\)
(via the continued-fraction fixed point), \(e\) (via the factorial
series), \(\sqrt{2}\) (via the Babylonian method), and \(\ln 2\) (via
the alternating harmonic series), all emerge from a single three-role
recursive grammar parameterized by a tholonic triple \((N, D, C)\). The
three roles, Negotiation (\(N\)), Definition (\(D\)), and Contribution
(\(C\)), are functionally irreducible: they cannot be collapsed to fewer
than three even when two positions share the same numerical seed value.

Paper 1's structural contribution is the \emph{grammar}: the five
branches are distinguished by initial seeds and traversal rules, but the
role assignments are invariant. The present paper asks a different
question: if we treat the seed values not as fixed canonical inputs but
as a space to be swept, what does the distribution of convergent
outcomes reveal about the five constants?

\subsection{The sweep question}\label{the-sweep-question}

The five algorithms each accept four parameters drawn from the tholonic
NDC model. When those parameters are drawn from the canonical seed set
\(\{1, 2, 3, 5\}\), the \(4^4 = 256\) combinations per algorithm yield
1,280 convergence runs. For each run we record whether the output
converges to a known constant. The result is a frequency table: how many
seed combinations send each algorithm to each constant?

That table turns out to have a striking structure. Every nonzero count
is an exact power of two, and those powers are \(2^8\), \(2^8\),
\(2^6\), \(2^4\), \(2^0\). This is not guaranteed by any symmetry of the
algorithms; it is an empirical result of the sweep. The present paper
documents, interprets, and extends that result.

\subsection{Relation to the existing series}\label{relation-to-the-existing-series}

This paper is a companion to Paper 1 and builds on its notation and
convergence proofs. Paper 3 {[}Mil26b{]} establishes the irreducibility
of three roles and is the foundational minimality result. Paper 6
{[}Mil26c{]} establishes that the qualitative character of the integers
1, 2, and 3 determines the role assignment. The present paper sits
closest to Paper 1: both concern the five constants under tholonic
recursion, but Paper 1 establishes the structural grammar while this
paper explores the combinatorial topology of the seed space.

\textbf{What this paper provides.} A full description of the
computational sweep methodology; the frequency table of convergences; a
proof that the counts are exact powers of two (given the seed set
\(\{1, 2, 3, 5\}\) and the five algorithms as specified); an argument
that the hierarchy is bound to the tholonic axis geometry of Paper 1;
the entropy interpretation; the divide-by-four cascade; the restoration
of \(\pi\) from the raw \(\pi/4\) output; and an argument that the
hierarchy reflects the algebraic standing of the constants.

\textbf{What this paper does not provide.} New convergence proofs for
the individual series (those are in Paper 1 and the classical
literature); any claim that the seed set \(\{1, 2, 3, 5\}\) is the only
set with this property; or a formal theorem characterizing all seed sets
that produce a power-of-two frequency table.

\textbf{Organization.} Section 2 defines the algorithms and the
four-parameter NDC seeding convention. Section 3 describes the sweep and
records the raw results. Section 4 establishes the power-of-two
structure and its binding to the tholonic axis geometry. Section 5
interprets the hierarchy through generative entropy. Section 6 examines
\(\varphi\) and \(\sqrt{2}\) as universal attractors. Section 7
reconstructs the complete cascade including \(\pi\) at the \(2^2\) slot.
Section 8 argues for \(\pi\) as archetype rather than instance. Section
9 states open questions. Section 10 concludes.

\begin{center}\rule{0.5\linewidth}{0.5pt}\end{center}

\section{Algorithms and Seed Parameterization}\label{algorithms-and-seed-parameterization}

\subsection{The tholonic NDC seeding model}\label{the-tholonic-ndc-seeding-model}

Each tholonic algorithm is seeded with four parameters drawn from the
NDC model:

\begin{longtable}[]{@{}
  >{\raggedright\arraybackslash}p{(\linewidth - 4\tabcolsep) * \real{0.4400}}
  >{\raggedright\arraybackslash}p{(\linewidth - 4\tabcolsep) * \real{0.3200}}
  >{\raggedright\arraybackslash}p{(\linewidth - 4\tabcolsep) * \real{0.2400}}@{}}
\toprule\noalign{}
\begin{minipage}[b]{\linewidth}\raggedright
Parameter
\end{minipage} & \begin{minipage}[b]{\linewidth}\raggedright
Symbol
\end{minipage} & \begin{minipage}[b]{\linewidth}\raggedright
Role
\end{minipage} \\
\midrule\noalign{}
\endhead
\bottomrule\noalign{}
\endlastfoot
Negotiation & \(N\) & Starting value; the state being iteratively
refined \\
Instantiation & \(I\) & Step size; drives the recursion depth \\
Definition & \(D\) & Limitation seed; bounds the recurrence \\
Contribution & \(C\) & Contribution seed; accumulates toward the
limit \\
\end{longtable}

The roles are functionally distinct. Two parameters may share the same
numerical value while playing structurally different roles (see Paper 3,
Lemma 4.2).

\subsection{The five algorithms}\label{the-five-algorithms}

Five recurrence relations implement the five algorithms:

\begin{longtable}[]{@{}
  >{\raggedright\arraybackslash}p{(\linewidth - 6\tabcolsep) * \real{0.1408}}
  >{\raggedright\arraybackslash}p{(\linewidth - 6\tabcolsep) * \real{0.1549}}
  >{\raggedright\arraybackslash}p{(\linewidth - 6\tabcolsep) * \real{0.4366}}
  >{\raggedright\arraybackslash}p{(\linewidth - 6\tabcolsep) * \real{0.2676}}@{}}
\toprule\noalign{}
\begin{minipage}[b]{\linewidth}\raggedright
Constant
\end{minipage} & \begin{minipage}[b]{\linewidth}\raggedright
Algorithm
\end{minipage} & \begin{minipage}[b]{\linewidth}\raggedright
Canonical Seeds \((N, I, D, C)\)
\end{minipage} & \begin{minipage}[b]{\linewidth}\raggedright
Convergence class
\end{minipage} \\
\midrule\noalign{}
\endhead
\bottomrule\noalign{}
\endlastfoot
\(\varphi\) & Continued fraction: \(N \leftarrow 1 + 1/N\) &
\((1, 2, 2, 3)\) & Class B (self-redefined) \\
\(\sqrt{2}\) & Babylonian: \(N \leftarrow (N + 2/N)/2\) &
\((1, 2, 2, 2)\) & Class B (self-redefined) \\
\(\ln 2\) & Alternating series: \(N \pm 1/(\text{count}+1)\) &
\((1, 2, 1, 2)\) & Class C (fixed) \\
\(e\) & Factorial series: \(N \leftarrow N + 1/C\);
\(C \leftarrow C \cdot \text{count}\) & \((1, 2, 1, 1)\) & Class B
(self-redefined) \\
\(\pi/4\) & Leibniz: \(N \leftarrow N - 1/D + 1/C\);
\(D, C \leftarrow D + I^2\), \(C + I^2\) & \((1, 2, 3, 5)\) & Class A
(advancing) \\
\end{longtable}

The convergence classification follows Paper 1, Section 4. Each run
iterates until successive outputs agree to six decimal places.

The seeds for \(\pi/4\), namely \((1, 2, 3, 5)\), are the opening four
terms of the Fibonacci sequence. This is not a coincidence imposed by
the framework; it emerges from the requirement that \(D\) and \(C\) must
be numerically distinct and that the step \(I^2 = 4\) is the unique
advance consistent with the Leibniz series's odd-denominator structure
(Paper 1, Section 5.5).

\begin{center}\rule{0.5\linewidth}{0.5pt}\end{center}

\section{The Sweep}\label{the-sweep}

\subsection{Design}\label{design}

The canonical seed set is \(S = \{1, 2, 3, 5\}\). We draw
\((N, I, D, C) \in S^4\), giving \(|S|^4 = 256\) input combinations.
Each combination is submitted to each of the five algorithms, yielding
\(256 \times 5 = 1{,}280\) total runs.

For each run we record the converged output and compare it against the
five known constants with a tolerance of \(\pm 5 \times 10^{-6}\). Runs
whose output does not match any of the five constants within this
tolerance are recorded as unlabeled convergences.

\subsection{Results}\label{results}

Of the 1,280 runs:

\begin{longtable}[]{@{}ll@{}}
\toprule\noalign{}
Outcome & Count \\
\midrule\noalign{}
\endhead
\bottomrule\noalign{}
\endlastfoot
Converged to a known constant & 593 (46.3\%) \\
Converged to an unlabeled value & 687 (53.7\%) \\
\end{longtable}

Breakdown by constant:

\begin{longtable}[]{@{}llll@{}}
\toprule\noalign{}
Constant & True value & Matching runs & Count as power of 2 \\
\midrule\noalign{}
\endhead
\bottomrule\noalign{}
\endlastfoot
\(\varphi\) & 1.618034\ldots{} & 256 & \(2^8\) \\
\(\sqrt{2}\) & 1.414214\ldots{} & 256 & \(2^8\) \\
\(\ln 2\) & 0.693147\ldots{} & 64 & \(2^6\) \\
\(e\) & 2.718282\ldots{} & 16 & \(2^4\) \\
\(\pi/4\) & 0.785398\ldots{} & 1 & \(2^0\) \\
\textbf{Total labeled} & & \textbf{593} & \\
\end{longtable}

The total labeled count, 593, is prime.

\begin{center}\rule{0.5\linewidth}{0.5pt}\end{center}

\section{The Power-of-Two Structure}\label{the-power-of-two-structure}

\subsection{Statement}\label{statement}

\textbf{Observation 4.1.} For the seed set \(S = \{1, 2, 3, 5\}\) and
the five algorithms as defined in Section 2.2, the number of seed
combinations in \(S^4\) that cause each algorithm to converge to its
associated constant is an exact power of two: \(256 = 2^8\) for
\(\varphi\) and \(\sqrt{2}\); \(64 = 2^6\) for \(\ln 2\); \(16 = 2^4\)
for \(e\); and \(1 = 2^0\) for \(\pi/4\).

\subsection{Why this is non-trivial}\label{why-this-is-non-trivial}

A combinatorial analysis of the seed space does not in general produce
powers of two. For arbitrary algorithms and arbitrary seed sets, the
count of ``successful'' inputs can be any integer from 0 to 256. The
fact that every nonzero count here is a power of two, and that those
powers are evenly spaced at intervals of 2 (with the apparent gap at
\(2^2\) filled when \(\pi\) is restored; see Section 7), is a structural
property of the interaction between the algorithms and the seed set
\(S\).

\subsection{Explanation for the universal cases}\label{explanation-for-the-universal-cases}

The counts for \(\varphi\) and \(\sqrt{2}\) follow immediately from the
mathematics of their algorithms:

\begin{itemize}
\item
  The continued-fraction iteration \(N \leftarrow 1 + 1/N\) converges to
  \(\varphi\) from any positive starting value. Since all elements of
  \(S\) are positive, all \(256\) combinations converge. The count is
  \(|S^4| = 2^8\).
\item
  The Babylonian iteration \(N \leftarrow (N + 2/N)/2\) converges to
  \(\sqrt{2}\) from any positive starting value by the AM-GM inequality.
  The same argument gives \(256 = 2^8\).
\end{itemize}

For \(\ln 2\), \(e\), and \(\pi/4\), the constraints are tighter. The
counts \(64 = 2^6\), \(16 = 2^4\), and \(1 = 2^0\) reflect the
increasing selectivity of those algorithms with respect to the seeding.
A complete combinatorial characterization is left as an open problem
(see Section 9, Conjecture 9.1).

\subsection{The exponents descend by two}\label{the-exponents-descend-by-two}

Reading the powers in convergence order:

\[
\varphi: 2^8 \quad \sqrt{2}: 2^8 \quad \ln 2: 2^6 \quad e: 2^4 \quad \pi/4: 2^0
\]

Each step is a factor of \(4 = 2^2\). The apparent jump from \(2^4\) to
\(2^0\) (skipping \(2^2\)) is resolved in Section 7.

\subsection{Binding to the tholonic axis geometry}\label{binding-to-the-tholonic-axis-geometry}

Paper 1 {[}Mil26a, Section 2.1{]} derives the five algorithms from a
geometric thologram in which six vertices carry successive powers of
two:

\begin{longtable}[]{@{}lll@{}}
\toprule\noalign{}
Vertex & Role position & Value \\
\midrule\noalign{}
\endhead
\bottomrule\noalign{}
\endlastfoot
Outer \(N\) & Negotiation & \(2^0 = 1\) \\
Outer \(D\) & Definition & \(2^1 = 2\) \\
Outer \(C\) & Contribution & \(2^2 = 4\) \\
Inner \(\mathrm{mid}_1\) & & \(2^3 = 8\) \\
Inner \(\mathrm{mid}_2\) & & \(2^4 = 16\) \\
Inner \(\mathrm{mid}_3\) & & \(2^5 = 32\) \\
\end{longtable}

Summing along the three directed axes yields totals that factor as \(7\)
times a characteristic multiplier. The axis multipliers are:

\begin{itemize}
\tightlist
\item
  Definition axis: \(5\)
\item
  Contribution axis: \(3\)
\item
  Instantiation axis: \(\mathrm{istep} = 2\)
\end{itemize}

The \(\pi/4\) branch alone advances externally with step
\(\Delta = \mathrm{istep}^2 = 2^2 = 4\), and its canonical seed triple
is \((N_0, D_0, C_0) = (1, 3, 5)\), where \(3\) and \(5\) are the
Contribution and Definition axis multipliers and \(1 = 2^0\) is the
generative unit.

The sweep seed set \(S = \{1, 2, 3, 5\}\) is therefore not an arbitrary
four-element subset of the positive integers. It is the
\textbf{axis-derived seed vocabulary} of the thologram: the outer \(N\)
vertex (\(2^0\)), the Instantiation axis multiplier
(\(\mathrm{istep} = 2\)), and the Contribution and Definition axis
multipliers (\(3\) and \(5\)). The outer \(C\) vertex (\(2^2 = 4\)) and
the three inner vertices (\(2^3, 2^4, 2^5\)) are excluded from \(S\),
but they enter indirectly through the combinatorial structure of the
sweep.

\textbf{The ceiling \(2^8\).} The sweep explores
\(|S|^4 = 4^4 = 256 = 2^8\) combinations. The cardinality \(|S| = 4\) is
not arbitrary: \(S\) has four elements because it contains exactly
\(2^0\), \(\mathrm{istep}\), \(m_C\), and \(m_D\), one token per
axis-derived role. Numerically, \(|S| = 4 = 2^2\), which also equals the
outer \(C\) vertex value. Whether this equality is structurally
necessary or a suggestive coincidence is not proved here; it is noted as
consistent with the axis geometry. What follows directly is that the
ceiling \(2^8\) is the four-fold product of \(|S|\), one factor per
swept parameter \((N, I, D, C)\), and that universal attractors
(\(\varphi\), \(\sqrt{2}\)) saturate it: every axis-vocabulary
assignment converges.

\textbf{The \(\div 4\) cascade rate.} Each step down the hierarchy
divides the path count by exactly \(4 = 2^2\). Paper 1 proves that
\(\Delta = \mathrm{istep}^2 = 4\) is the unique Class A advance step for
the \(\pi/4\) branch, derived from the Instantiation axis multiplier.
The cascade rate of the seed-space hierarchy is numerically the same
\(2^2\) that drives parameter advancement in the \(\pi/4\) algorithm.
Whether this is a structural correspondence or a causal derivation is
not proved here: the claim is that the two quantities are equal and that
this equality is consistent with the axis geometry, not that one forces
the other by a theorem already in hand.

\textbf{The even-exponent ladder.} The convergence exponents form a
gapless sequence of even integers:

\[2^0, \quad 2^2, \quad 2^4, \quad 2^6, \quad 2^8.\]

The thologram vertex ladder runs \(2^0\) through \(2^5\) in unit
exponent steps. The hierarchy ladder extends further, to \(2^8\), so it
is not a subset of the vertex ladder. What it shares with the vertex
ladder is the binary base; what distinguishes it is that only even
exponents appear, spaced by \(\mathrm{istep} = 2\). Each step up the
hierarchy multiplies the count by \(\mathrm{istep}^2 = 4\). The
Instantiation axis multiplier thus sets the exponent step of the
hierarchy, even though it does not limit the hierarchy's range to that
of the vertex ladder itself. Whether this exponent-step correspondence
is a theorem or a structural echo is left to Conjecture 9.5.

\textbf{The singular origin.} The unique \(\pi/4\) path uses
\((N, I, D, C) = (1, 2, 3, 5)\): the full axis vocabulary in order.
Section 5.4 discusses the Fibonacci structure of this combination; from
the axis-geometry perspective, it is simply the complete set of
thologram-derived tokens assembled in a single convergent run, with zero
generative entropy. No other combination of \(S\)-values in any
algorithm reaches \(\pi/4\).

\textbf{What is proved vs.~what is argued.} The following bindings are
established in Paper 1: the vertex values \(2^0, 2^1, 2^2\); the axis
multipliers \(5, 3, 2\); the step \(\Delta = \mathrm{istep}^2 = 4\); and
the seed triple \((1, 3, 5)\) for the \(\pi/4\) branch. The following
bindings are argued here on the basis of the sweep data and structural
correspondence, but not yet proved as theorems: that the even-exponent
ladder \(0, 2, 4, 6, 8\) is the unique consequence of axis geometry;
that the intermediate counts \(2^6\) and \(2^4\) for \(\ln 2\) and \(e\)
are forced by the axis structure rather than by algorithm-specific
combinatorics alone; and that no non-axis-derived seed set reproduces
the same gapless hierarchy. Conjecture 9.5 (Section 9) states the last
of these formally.

The power-of-two pattern is therefore not merely an empirical curiosity
about five classical series. It is the \textbf{binary vertex geometry of
the thologram expressing itself combinatorially} when the axis-derived
seed vocabulary is swept across the five tholonic branches. The
hierarchy is tholonic not only in its algorithms and roles, but in the
powers of two that count its convergent paths.

\begin{center}\rule{0.5\linewidth}{0.5pt}\end{center}

\subsection{\texorpdfstring{5. Generative Entropy and the Singularity of
\(\pi/4\)}{5. Generative Entropy and the Singularity of \textbackslash pi/4}}\label{generative-entropy-and-the-singularity-of-pi4}

\subsection{Definition}\label{definition}

Define the \emph{generative entropy} of a constant \(\kappa\) (with
respect to the sweep) as the number of bits required to specify which
seed combination produces \(\kappa\):

\[H(\kappa) = \log_2(\text{number of seed combinations producing } \kappa).
\]

\begin{longtable}[]{@{}lll@{}}
\toprule\noalign{}
Constant & Paths & \(H(\kappa)\) (bits) \\
\midrule\noalign{}
\endhead
\bottomrule\noalign{}
\endlastfoot
\(\pi/4\) & 1 & \(\log_2(1) = 0\) \\
\(e\) & 16 & \(\log_2(16) = 4\) \\
\(\ln 2\) & 64 & \(\log_2(64) = 6\) \\
\(\varphi\) & 256 & \(\log_2(256) = 8\) \\
\(\sqrt{2}\) & 256 & \(\log_2(256) = 8\) \\
\end{longtable}

\subsubsection{\texorpdfstring{5.2 \(\pi/4\) has zero generative
entropy}{5.2 \textbackslash pi/4 has zero generative entropy}}\label{pi4-has-zero-generative-entropy}

There is exactly one seed combination that produces \(\pi/4\):
\((N, I, D, C) = (1, 2, 3, 5)\) applied to the Leibniz algorithm. No
information is needed to select it, because no choice exists. The path
is uniquely determined.

This is a striking inversion of how \(\pi\) is ordinarily characterized.
\(\pi\) is typically regarded as one of the most
\emph{representationally complex} constants: transcendental, infinite,
non-repeating, without closed-form algebraic expression. Yet within this
seed-space framework it is the most \emph{ordered}, the most
constrained, the most ``prime-like'' in the sense of being irreducible
to simpler combinatorial constituents.

\subsection{Two distinct kinds of complexity}\label{two-distinct-kinds-of-complexity}

The sweep reveals a distinction between two fundamentally different
senses of complexity:

\begin{longtable}[]{@{}
  >{\raggedright\arraybackslash}p{(\linewidth - 4\tabcolsep) * \real{0.2619}}
  >{\raggedright\arraybackslash}p{(\linewidth - 4\tabcolsep) * \real{0.2143}}
  >{\raggedright\arraybackslash}p{(\linewidth - 4\tabcolsep) * \real{0.5238}}@{}}
\toprule\noalign{}
\begin{minipage}[b]{\linewidth}\raggedright
Dimension
\end{minipage} & \begin{minipage}[b]{\linewidth}\raggedright
\(\pi/4\)
\end{minipage} & \begin{minipage}[b]{\linewidth}\raggedright
\(\varphi\), \(\sqrt{2}\)
\end{minipage} \\
\midrule\noalign{}
\endhead
\bottomrule\noalign{}
\endlastfoot
\textbf{Representational complexity} & High: transcendental, infinite,
non-repeating & Low: roots of minimal quadratics \\
\textbf{Generative entropy} & Zero: one unique path & Maximum: all 256
paths converge \\
\end{longtable}

\(\pi/4\) is hard to \emph{describe} but uniquely \emph{located} in the
seed space. \(\varphi\) and \(\sqrt{2}\) are easy to \emph{describe}
algebraically but require more bits to specify \emph{which} path was
taken to reach them.

\subsection{Uniqueness of the Fibonacci seeds}\label{uniqueness-of-the-fibonacci-seeds}

The one combination that produces \(\pi/4\) is \((1, 2, 3, 5)\), the
first four terms of the Fibonacci sequence. This is the same sequence
whose ratio converges to \(\varphi\). The seeds that uniquely unlock
\(\pi/4\) are simultaneously the seeds that encode the golden ratio's
generative structure. This suggests a deep connection at the level of
the generative apparatus, not merely at the level of the constants
themselves.

\begin{center}\rule{0.5\linewidth}{0.5pt}\end{center}

\subsection{\texorpdfstring{6. \(\varphi\) and \(\sqrt{2}\) as Universal
Attractors}{6. \textbackslash varphi and \textbackslash sqrt\{2\} as Universal Attractors}}\label{varphi-and-sqrt2-as-universal-attractors}

\subsection{Shared properties}\label{shared-properties}

\(\varphi\) and \(\sqrt{2}\) are the only constants in the hierarchy
that occupy the same level (\(2^8\)). Every other level has exactly one
occupant. Their shared universality reflects several concurrent
structural facts:

\textbf{Algebraic minimality.} Both are roots of the simplest possible
quadratics: - \(\varphi\) satisfies \(x^2 - x - 1 = 0\) - \(\sqrt{2}\)
satisfies \(x^2 - 2 = 0\)

These are the two minimal quadratic irrationals: one more step above the
integers than any quadratic can be.

\textbf{Metallic-mean adjacency.} Both constants also sit on adjacent
rungs of the \emph{metallic-mean} ladder. The metallic mean \(M_n\) is
the positive root of \(x^2 - nx - 1 = 0\), equivalently
\(M_n = (n + \sqrt{n^2 + 4})/2\). For \(n = 1\), \(M_1 = \varphi\) (the
golden mean). For \(n = 2\), \(M_2 = 1 + \sqrt{2}\) (the silver mean).
Substituting \(y = x - 1\) into the silver-mean equation
\((y+1)^2 - 2(y+1) - 1 = 0\) yields \(y^2 - 2 = 0\), so
\(y = \sqrt{2}\). Thus \(\sqrt{2} = M_2 - 1\): not itself a metallic
mean, but the value one unit below the silver mean on the same ladder.
\(\varphi\) occupies \(M_1\) directly; \(\sqrt{2}\) occupies
\(M_2 - 1\). The sisterhood is therefore not only algorithmic
universality but also adjacency in the metallic-mean family at the first
two indices.

\textbf{Self-referential fixed-point structure.} Both algorithms are
self-referential in a deep sense: - \(\varphi = 1 + 1/\varphi\) (the
continued fraction contains itself) -
\(\sqrt{2} = (\sqrt{2} + 2/\sqrt{2})/2\) (the Babylonian average
reconstructs itself)

This structural self-reference is what makes the algorithms converge
from any positive starting point.

\textbf{Geometric ubiquity.} \(\sqrt{2}\) is the diagonal of the unit
square; \(\varphi\) is the diagonal-to-side ratio of the regular
pentagon. Both appear throughout natural proportion, tiling, and growth
patterns. Their geometric generality parallels their algorithmic
universality.

\subsection{A natural polarity}\label{a-natural-polarity}

The hierarchy exhibits a natural polarity between the universal
attractors and the singular origin:

\begin{itemize}
\tightlist
\item
  \(\varphi\) and \(\sqrt{2}\) are \textbf{universal receivers}: present
  everywhere, reachable from any starting point, requiring no special
  key.
\item
  \(\pi/4\) is the \textbf{singular source}: present in exactly one
  place, requiring the precise Fibonacci key \((1, 2, 3, 5)\) to unlock
  it.
\item
  \(e\) and \(\ln 2\) are \textbf{intermediaries}: each a factor of 4
  more accessible than the last, forming the cascade connecting source
  to receiver.
\end{itemize}

\begin{center}\rule{0.5\linewidth}{0.5pt}\end{center}

\subsection{\texorpdfstring{7. The Complete Cascade: Restoring \(\pi\)
at the \(2^2\)
Slot}{7. The Complete Cascade: Restoring \textbackslash pi at the 2\^{}2 Slot}}\label{the-complete-cascade-restoring-pi-at-the-22-slot}

\subsection{The divide-by-four cascade}\label{the-divide-by-four-cascade}

The seed-space counts form a cascade that divides by 4 at each step:

\[
256 \xrightarrow{\div 4} 64 \xrightarrow{\div 4} 16 \xrightarrow{\div 4} 4 \xrightarrow{\div 4} 1
\]

The raw sweep records \(1\) at the final position (for \(\pi/4\)). The
predicted intermediate position of \(4\) appears unoccupied.

\subsubsection{\texorpdfstring{7.2 The tholonic Class A algorithm
outputs \(\pi/4\), not
\(\pi\)}{7.2 The tholonic Class A algorithm outputs \textbackslash pi/4, not \textbackslash pi}}\label{the-tholonic-class-a-algorithm-outputs-pi4-not-pi}

This apparent gap is resolved by attending to what the tholonic
algorithm actually computes. The Class A branch for \(\pi/4\) is derived
from the axis geometry of the thologram: the seeds
\((N_0, D_0, C_0) = (1, 3, 5)\) come from the axis multipliers and the
generative unit, and the advance step \(\Delta = \mathrm{istep}^2 = 4\)
is derived from the Instantiation axis (Paper 1, Section 2.1). The
algorithm was not constructed from the Leibniz series; the Leibniz
series is what the algorithm reduces to mathematically, as a consequence
of its structure (Paper 1, Proposition 5.5).

The raw output of this algorithm is \(\pi/4\). Recovering \(\pi\)
requires an explicit \(\times 4\) scaling step that is external to the
algorithm itself. When that scaling is applied, the one seed combination
producing \(\pi/4\) corresponds to four combinations producing \(\pi\)
(one per element of the \(\times 4\) factor). The \(2^2 = 4\) slot is
not empty; it is occupied by \(\pi\).

This is not a rescaling convention borrowed from the Leibniz literature.
It is an observation about the tholonic algorithm's native output: the
algorithm converges to \(\pi/4\) as a structural consequence of its
Class A geometry, and \(\pi\) itself sits one scaling step above that
native value.

\subsection{The complete hierarchy}\label{the-complete-hierarchy}

\begin{longtable}[]{@{}lll@{}}
\toprule\noalign{}
Constant & Combinations & Power of 2 \\
\midrule\noalign{}
\endhead
\bottomrule\noalign{}
\endlastfoot
\(\varphi\) & 256 & \(2^8\) \\
\(\sqrt{2}\) & 256 & \(2^8\) \\
\(\ln 2\) & 64 & \(2^6\) \\
\(e\) & 16 & \(2^4\) \\
\(\pi\) & 4 & \(2^2\) \\
\(\pi/4\) & 1 & \(2^0\) \\
\end{longtable}

Every even power of 2 from \(2^0\) to \(2^8\) is occupied. The sequence
is complete and gapless:

\[
2^0 \quad 2^2 \quad 2^4 \quad 2^6 \quad 2^8 \quad 2^8
\] \[
\pi/4 \quad \pi \quad e \quad \ln 2 \quad \sqrt{2} \quad \varphi
\]

\subsection{The four-stage emanation structure}\label{the-four-stage-emanation-structure}

Grouping by accessibility:

\begin{longtable}[]{@{}llll@{}}
\toprule\noalign{}
Stage & Constants & Combinations & Multiplier from previous stage \\
\midrule\noalign{}
\endhead
\bottomrule\noalign{}
\endlastfoot
0 --- Singular & \(\pi/4\) & 1 & (origin) \\
1 --- Rare & \(e\) & 16 & \(\times 16\) \\
2 --- Common & \(\ln 2\) & 64 & \(\times 4\) \\
3 --- Universal & \(\varphi\), \(\sqrt{2}\) & 256 & \(\times 4\) \\
\end{longtable}

Each outward step multiplies accessibility by 4. \(\pi\) occupies the
slot between Stage 0 (\(\pi/4\)) and Stage 1 (\(e\)), at \(\times 4\)
from the origin.

\begin{center}\rule{0.5\linewidth}{0.5pt}\end{center}

\subsection{\texorpdfstring{8. \(\pi\) as Archetype Rather Than
Instance}{8. \textbackslash pi as Archetype Rather Than Instance}}\label{pi-as-archetype-rather-than-instance}

\subsection{The apparent paradox}\label{the-apparent-paradox}

\(\varphi\) and \(\sqrt{2}\) are the \emph{most common} constants in
this framework, appearing in every possible seed combination. Yet
\(\pi\) is fundamental to far more of mathematics and physics than
either of them: it appears in the normal distribution, the Fourier
transform, the Basel problem (\(1 + 1/4 + 1/9 + \cdots = \pi^2/6\)),
wave equations, and the geometry of all circles. How can the rarest
value in the seed space be the most fundamental value in mathematics?

\subsection{The archetype-instance distinction}\label{the-archetype-instance-distinction}

The resolution lies in the distinction between an \emph{archetype} and
an \emph{instance}:

\begin{itemize}
\tightlist
\item
  An \textbf{instance} is common and reproducible. It appears \emph{in}
  particular things. It is reachable from many starting points.
\item
  An \textbf{archetype} is singular and irreducible. It is the pattern
  that generates instances rather than being one itself. It is reached
  only by a specific path.
\end{itemize}

By this distinction the hierarchy maps precisely:

\begin{longtable}[]{@{}
  >{\raggedright\arraybackslash}p{(\linewidth - 2\tabcolsep) * \real{0.6250}}
  >{\raggedright\arraybackslash}p{(\linewidth - 2\tabcolsep) * \real{0.3750}}@{}}
\toprule\noalign{}
\begin{minipage}[b]{\linewidth}\raggedright
Constant
\end{minipage} & \begin{minipage}[b]{\linewidth}\raggedright
Role
\end{minipage} \\
\midrule\noalign{}
\endhead
\bottomrule\noalign{}
\endlastfoot
\(\varphi\), \(\sqrt{2}\) & Instances: universal attractors, appear
everywhere \\
\(\ln 2\), \(e\) & Mediators: bridge between archetype and instance \\
\(\pi\) & Archetype: singular, irreducible, the source \\
\end{longtable}

\subsubsection{\texorpdfstring{8.3 Consistency with \(\pi\)'s
mathematical
behavior}{8.3 Consistency with \textbackslash pi's mathematical behavior}}\label{consistency-with-pis-mathematical-behavior}

This interpretation is consistent with how \(\pi\) actually behaves
across mathematics. \(\varphi\) and \(\sqrt{2}\) appear in
\emph{specific structures}: this spiral, this rectangle, this triangle.
\(\pi\) appears in the \emph{laws governing all structures}: in the
relationship between circumference and diameter, in the structure of the
normal distribution, in the completeness of the Fourier basis, in
quantum mechanics. It shows up not in particular instances but in the
framework within which all instances are measured.

An archetype does not need to be common. Its rarity is its power. If
\(\pi\) appeared in all 256 combinations it would be just another
attractor. The fact that it requires the one specific Fibonacci key
\((1, 2, 3, 5)\) suggests it is not arrived at by exploration or
accident: every other constant can be stumbled upon from any starting
point, but \(\pi\) must be sought.

\begin{center}\rule{0.5\linewidth}{0.5pt}\end{center}

\section{Open Questions}\label{open-questions}

The following conjectures are stated informally; resolving them would
elevate the empirical observations of this paper to theorems.

\textbf{Conjecture 9.1 (Power-of-two characterization).} The seed set
\(\{1, 2, 3, 5\}\) is not the only set for which the frequency table of
convergences is composed entirely of exact powers of two. Characterize
all finite seed sets \(S \subset \mathbb{R}_{>0}\) with this property.

\textbf{Conjecture 9.2 (Gapless even-power sequences).} For the five
algorithms as defined, any seed set satisfying Conjecture 9.1 also
produces a frequency table in which the exponents form a gapless
sequence of even integers from 0 to \(2\log_2|S^4|\). Prove or disprove.

\textbf{Conjecture 9.3 (Fibonacci necessity for \(\pi/4\)).} Among all
seed sets satisfying Conjecture 9.1, the unique combination producing
\(\pi/4\) always consists of consecutive Fibonacci numbers. Prove or
find a counterexample.

\textbf{Conjecture 9.4 (Unlabeled convergences).} The 687 runs
converging to unlabeled values in Section 3.2 also exhibit a
power-of-two frequency table when grouped by distinct convergent value.
Verify computationally and characterize the unlabeled constants.

\textbf{Conjecture 9.5 (Axis-geometry necessity).} A finite seed set
\(S\) produces the gapless even-exponent hierarchy
\(2^0, 2^2, \ldots, 2^{2\log_2|S|}\) under the five Paper 1 algorithms
if and only if \(S\) is the axis-derived vocabulary
\(\{2^0, \mathrm{istep}, m_C, m_D\} = \{1, 2, 3, 5\}\), where \(m_C\)
and \(m_D\) are the Contribution and Definition axis multipliers. Prove
or find a counterexample.

\begin{center}\rule{0.5\linewidth}{0.5pt}\end{center}

\section{Conclusion}\label{conclusion}

A computational sweep of all \(4^4 = 256\) seed combinations from
\(\{1, 2, 3, 5\}\) across five tholonic algorithms reveals that the
number of combinations converging to each classical constant is an exact
power of two. The powers are \(2^8\), \(2^8\), \(2^6\), \(2^4\), and
\(2^0\), and when the Leibniz algorithm's raw output \(\pi/4\) is scaled
by 4 to restore \(\pi\), the slot \(2^2 = 4\) is filled and the sequence
of even powers from \(2^0\) to \(2^8\) is complete and gapless.

This cascade divides by 4 at each step and admits a clean entropy
interpretation: \(\pi/4\) has zero generative entropy (one path, no
choice), while \(\varphi\) and \(\sqrt{2}\) have maximum generative
entropy (all paths lead there). The intermediate constants \(e\) and
\(\ln 2\) form a bridge, each a factor of 4 more accessible than the
last.

The hierarchy mirrors the algebraic standing of the constants: quadratic
irrationals at maximum accessibility, transcendentals at minimum
accessibility. Section 4.5 further argues that the power-of-two
structure is bound to the tholonic axis geometry: the sweep seed set
\(\{1, 2, 3, 5\}\) is axis-derived, the cascade rate \(\div 4\) equals
\(\mathrm{istep}^2\), and the even-exponent ladder propagates the binary
vertex structure through the Instantiation axis. We argue this is not
coincidence but reflects what the constants structurally are within the
tholonic framework, and that \(\pi\)'s position as the singular source
in the seed space is consistent with its role as an archetype rather
than an instance in the broader structure of mathematics.

\begin{center}\rule{0.5\linewidth}{0.5pt}\end{center}

\section{References}\label{references}

{[}Fib02{]} Fibonacci (Leonardo of Pisa). \emph{Liber Abaci.} 1202.
(Standard modern translation: Sigler, L. E. \emph{Fibonacci's Liber
Abaci.} Springer, 2002.)

{[}Her97{]} Herstein, I. N. \emph{Abstract Algebra.} 3rd ed.~Wiley,
1997.

{[}Knu97{]} Knuth, D. E. \emph{The Art of Computer Programming, Vol. 2:
Seminumerical Algorithms.} 3rd ed.~Addison-Wesley, 1997.

{[}Lei82{]} Leibniz, G. W. ``Nova methodus pro maximis et minimis.''
\emph{Acta Eruditorum,} 1684. (Cited for the series identity
\(\pi/4 = 1 - 1/3 + 1/5 - \cdots\), proved in {[}Roy11{]}.)

{[}Mil26a{]} Milton, J. W. ``Emergence of Classical Constants from a
Minimal Recursive Triadic Framework.'' Clarity Coalition preprint, May
2026. {[}Paper 1 of this series.{]}

{[}Mil26b{]} Milton, J. W. ``A Minimal Recursive Triadic Framework:
Irreducibility and Structural Necessity.'' Clarity Coalition preprint,
April 2026. {[}Paper 3 of this series.{]}

{[}Mil26c{]} Milton, J. W. ``The Qualitative Nature of One, Two, and
Three: Structural Role Assignment in Minimal Recursive Systems.''
Clarity Coalition preprint, April 2026. {[}Paper 6 of this series.{]}

{[}Roy11{]} Roy, R. \emph{Sources in the Development of Mathematics:
Series and Products from the Fifteenth to the Twenty-First Century.}
Cambridge University Press, 2011.

{[}Ste92{]} Stewart, I. \emph{Another Fine Math You've Got Me Into.} W.
H. Freeman, 1992.

\end{document}
